%%%%%%%%%%%%%%%%%%%%%%%%%%%%%%%%%%%%%%%%%%%%%%%%%%%%%%%%%%%%%%%%%%%%%%%%%%%%%%%%
%%  VORLAGE FÜR FREIWILLIGE WEITERE EINLAGEN
%%  Datei kopieren (z.B. 7_datenblatt.tex), in Endabgabe einbinden:
%%      \input{einlagen/7_datenblatt}
%%  "Sie können eine beliebige Anzahl (>0) weiterer Einlagen zu Ihrem
%%   Schwerpunktthema einreichen. Die Struktur ist die gleiche wie bei den
%%   sechs Pflichteinlagen."  Beispiele laut Leitfaden: Grafik, Buchkapitel,
%%   Handout, Gesprächsnotiz, Datenblatt, selbst angefertigtes Objekt,
%%   Video-/Tonaufzeichnung, klassisches Protokoll.
%%%%%%%%%%%%%%%%%%%%%%%%%%%%%%%%%%%%%%%%%%%%%%%%%%%%%%%%%%%%%%%%%%%%%%%%%%%%%%%%
\section{Einlage: \platzhalter{Titel der Einlage}}
\label{sec:zusatz1}

\subsection{Beschreibung}
\begin{hinweise}
Um welche Art von Einlage handelt es sich? Woher stammt diese Einlage, wie bin ich darauf aufmerksam
geworden? Welche sind die wesentlichen Aussagen dieser Einlage (eigene Worte)? ($1/4$ bis max.\ $1/2$ Seite)
\end{hinweise}
\begin{umfang}{Zusatzeinlage: Beschreibung}{0.25}{0.5}
\platzhalter{Beschreibung}
\end{umfang}

%% Nur falls physikalischer Inhalt (dann 1–3 Seiten):
%\subsection{Diskussion der Physik}
%\begin{umfang}{Zusatzeinlage: Diskussion der Physik}{1}{3}
%...
%\end{umfang}

\subsection{Reflexion}
\begin{hinweise}
Wie habe ich mit dieser Einlage gearbeitet? Warum habe ich sie ausgewählt? Bezug zum Schwerpunktthema?
Wichtigste Erkenntnisse jenseits des wissenschaftlichen Inhalts? Bekanntes vertieft, neu gelernt, neue Fragen?
Einfluss auf den weiteren Lernprozess -- welche Konsequenzen? Was sagen Kommilitoninnen/Kommilitonen dazu?
Unstimmigkeiten? Überraschende Erkenntnis? (1--2 Seiten)
\end{hinweise}
\begin{umfang}{Zusatzeinlage: Reflexion}{1}{2}
\platzhalter{Reflexion}
\end{umfang}
